\section{Conclusions}
\label{sec:conclusion}

We have developed a physics-based model for the acoustic `thump test''
used to assess watermelon ripeness.  The principal findings are:

\begin{enumerate}
  \item The watermelon is modelled as a fluid-filled viscoelastic shell
    using an equivalent-sphere approximation: the oblate geometry is
    captured through $R_{\mathrm{eq}} = (a^2 c)^{1/3}$, and the
    eigenfrequency is derived from classical thin-shell theory for a
    sphere of that radius.  The dominant tap-tone corresponds to the
    $n = 2$ flexural mode.

  \item The squared frequency $f_2^2$ is exactly linear in the rind
    elastic modulus $E_{\mathrm{rind}}$ when turgor pressure is
    negligible.  This linearity enables a closed-form inversion
    (Eq.~\ref{eq:inversion}) from a single tap-tone measurement to the
    rind modulus.  The condition number of this inversion is unity with
    respect to~$f^2$ and approximately~2 with respect to~$f$.

  \item Sobol sensitivity analysis confirms that
    $E_{\mathrm{rind}}$ dominates the frequency response
    ($S_T = 0.54 \pm 0.05$), with fruit size as the second most influential
    parameter.

  \item The dimensionless group $\Pi_{\mathrm{ripe}}$
    (Eq.~\ref{eq:Pi_ripe}) is a geometric invariant: $E_{\mathrm{rind}}$
    cancels by construction at fixed geometry, making it independent of
    the elastic modulus.
    Its practical value lies in collapsing different cultivar geometries
    onto a single calibration constant
    ($\Pi_{\mathrm{ripe}} = 0.062 \pm 0.002$ across four cultivars).

  \item Model predictions are benchmarked against published resonance
    measurements on watermelons \cite{Yamamoto1980} and are consistent
    with frequency--stiffness correlations reported for other fruits
    \cite{DeBelieetal2000, ChenDeBaerdemaeker1993}.
\end{enumerate}

The equivalent-sphere approximation is adequate for the inversion
application because the goal is to recover~$E_{\mathrm{rind}}$ from
$f_2$ at known geometry; since $a$, $c$, and~$h$ are measured
independently, $R_{\mathrm{eq}}$ is determined and the inversion is
well-posed.  However, the model has zero sensitivity to aspect ratio at
constant~$R_{\mathrm{eq}}$: it cannot distinguish a sphere from an
oblate spheroid of equal volume.  A Ritz variational model with explicit
$\zeta$-dependent basis functions would be needed to capture genuine
curvature effects and could improve accuracy for strongly oblate
cultivars ($\zeta < 0.6$).

The framework generalises directly from prior work on biological
cavities, confirming that classical shell theory provides a unified
language for fluid-filled biological structures---whether the application
is clinical diagnostics or selecting the perfect watermelon.

Future work will pursue experimental validation on instrumented
watermelons, extension of the model to include explicit oblate curvature
corrections via a Ritz variational formulation, application to other
fruits (cantaloupe, coconut, pumpkin), and development of a
smartphone-based measurement tool.
